\documentclass[10pt,a4paper]{amsart}
\usepackage[margin=19mm]{geometry}
\usepackage{amsmath,amssymb,mathtools}
\usepackage[scr=rsfs,cal=euler]{mathalfa}
\usepackage{xfrac}
\usepackage[unicode,hidelinks,bookmarks=false]{hyperref}
\newcommand{\ie}{i.e.\ }
\newcommand{\cf}{cf.\ }
\newcommand{\resp}{respectively\ }
\newcommand{\Subsection}{\subsection}
\providecommand{\nobreakdash}{\nobreak\hbox{-}}
\setlength{\emergencystretch}{2em}
\multlinegap0pt
\usepackage{bm}
\usepackage{bbm}
\usepackage{dsfont}
\usepackage{xspace}
\usepackage{cancel}
\usepackage{mathtools}
\usepackage{amsthm}
\makeatletter
\@ifundefined{theorem}{\newtheorem{theorem}{Theorem}}{}
\@ifundefined{proposition}{\newtheorem{proposition}[theorem]{Proposition}}{}
\makeatother
\newcommand{\R}{\mathbb{R}}
\newcommand{\Sig}{\mathcal{A}}
\newcommand{\norm}[1]{\left\|#1\right\|}
\title[From Gradient Descent to Potential Theory: A Geometric Dicti]{From Gradient Descent to Potential Theory: A Geometric Dictionary for Binary Classification}
\author{Catalin Vasii}
\date{Source: arXiv:2607.00988v2 (CC BY 4.0)}
\begin{document}
\maketitle

\noindent\emph{Source and licence.} From Gradient Descent to Potential Theory: A Geometric Dictionary for Binary Classification. Version arXiv:2607.00988v2 (CC BY 4.0), distributed under the Creative Commons Attribution 4.0 International licence, as recorded in the arXiv licence metadata for this exact version and shown on the abstract page https://arxiv.org/abs/2607.00988v2. Full rights notice: licenses/arxiv-2607.00988-CC-BY-4.0.txt.\par\smallskip

\noindent\textit{Editorial scope.} Counted unit: the Theorem whose source label is \texttt{thm:flat} in the article (source0.tex lines 1084--1090, source label \texttt{thm:flat}), delivered together with the article's own complete original proof of it (source0.tex lines 1092--1113, one \texttt{proof} environment of 22 lines). No mathematics is written, repaired or re-derived: statement and proof are the article's own text line for line. Required context retained uncounted: eq:curvature (lines 990--994); prop:flat-criterion (lines 998--1008). Every definition, display and label of those lines is retained unchanged. Omitted from this package and not redistributed: 1--989, 995--997, 1009--1083, 1091--1091, 1114--2649. Nothing inside a retained excerpt is omitted or altered: every retained body is delivered exactly as the article's lines are written, so no omission receipt of any kind accompanies this package. Citations: the retained excerpts carry no citation command at all, so no bibliography entry of the article is transcribed and no reference is redistributed. Reference adaptation: none was needed and none was made; every reference inside the delivered unit resolves inside it, so no unresolved reference or citation is produced. Printed slips kept literal: no printed word, symbol, hypothesis or number of the counted statement or of its proof was altered. Layout and class: the article's own class is \texttt{article} with packages including fontenc, inputenc, mathpazo, amsmath, amssymb, amsthm, geometry, hyperref, booktabs, array, tabularx, microtype, enumitem, xcolor; the wrapper is the lane's pinned \texttt{amsart} editorial wrapper at 10pt on A4, which restates the theorem-like environments the retained text needs and adds the article's own macro definitions that the retained text needs (\texttt{\textbackslash R}, \texttt{\textbackslash Sig}, \texttt{\textbackslash norm}), with extra wrapper packages (bm, bbm, dsfont, xspace, cancel, mathtools). The delivered numbering is the unit's own. Assets: no retained span contains a figure environment, a graphics-inclusion command, a PDF-page inclusion, a table or a TikZ picture; no image file is copied into this package, the package declares no asset dependency and no image file is redistributed with it. Licence: version arXiv:2607.00988v2 of this article is distributed under the Creative Commons Attribution 4.0 International licence, as recorded in the arXiv licence metadata for this exact version and shown on the abstract page https://arxiv.org/abs/2607.00988v2. A full rights notice is delivered at licenses/arxiv-2607.00988-CC-BY-4.0.txt. Characters: every retained excerpt is pure ASCII, so the delivered engine, XeLaTeX with its default Latin Modern text font, reports no missing character.
\begin{equation}\label{eq:curvature}
  \boxed{F_{\mu\nu}
  \;=\; \frac{1}{\norm{s}^2}
        \Bigl(a_\mu\,(a_\nu^\perp)^T - a_\nu\,(a_\mu^\perp)^T\Bigr).}
\end{equation}

\begin{proposition}\label{prop:flat-criterion}
Write $a_\mu = a_\mu^\perp + \alpha_\mu s$ with $a_\mu^\perp \in s^\perp$. The
induced connection \eqref{eq:curvature} is flat if and only if, at every
$x \in M$ and for all $\mu,\nu$, \emph{both}
\[
  \text{(a)}\quad a_\mu^\perp \wedge a_\nu^\perp = 0,
  \qquad\qquad
  \text{(b)}\quad \alpha_\mu\, a_\nu^\perp = \alpha_\nu\, a_\mu^\perp .
\]
No condition on the topology of $M$ is required.
\end{proposition}

\begin{theorem}[Existence of flat solutions]\label{thm:flat}
Let $M$ be a smooth manifold and $\mathcal{D} = \{(x_i, y_i)\}$ a finite binary
classification dataset. Then there exists a section $s: M \to \R^2
\hookrightarrow \R^m$ of constant norm and a flat $O(2)$-connection $\nabla$ on
$E = M \times \R^m$ such that $[\nabla, s]$ solves the classification problem
with readout $\varphi = (1,0,\ldots,0)$.
\end{theorem}

\begin{proof}
By a standard partition of unity argument, there exists a smooth function
$f: M \to [-1,1]$ with $f(x_i) = y_i$ for all $i$. Define
\[
  \theta(x) = \frac{\pi}{2}(1 - f(x)) \;\in\; [0,\pi],
\]
so that $f(x_i) = +1 \Rightarrow \theta(x_i) = 0$ and $f(x_i) = -1 \Rightarrow
\theta(x_i) = \pi$. Set
\[
  s(x) = r\bigl(\cos\theta(x),\;\sin\theta(x),\;0,\ldots,0\bigr)^T \in \R^m
\]
for any $r > 0$. The section takes values in a fixed 2-dimensional subspace, so
by Proposition~\ref{prop:flat-criterion} the induced connection is flat. The
horizontality condition gives $\Sig = -d\theta \cdot J$ where $J$ acts on the
first two components. The curvature is
\[
  F = d\Sig + \Sig\wedge\Sig = -d^2\theta\cdot J + (d\theta\wedge d\theta)\cdot J^2 = 0,
\]
since $d^2 = 0$ and $d\theta\wedge d\theta = 0$ for any 1-form. The classifier is
$\varphi(s(x)) = r\cos\theta(x)$, which has the correct sign at every data point
since $\cos 0 = +1$ and $\cos\pi = -1$.
\end{proof}

\end{document}
